\documentclass[10pt,a4paper]{article}
\usepackage{amsmath,amssymb,amsthm,graphicx}
% Card-local minimal preamble (journal unit sigma-2024-003-optimal-transport-and-generalized-ricci-flow, P2-ahl16 precedent).
%% Verbatim from the pinned source: the author theorem environments and macro block,
%% so every macro and environment of the counted spans is available.
% Omitted and why: the documentclass line (the journal class sigma is not on the host;
%% the assembler supplies the class) and page-geometry/font blocks (they fight the
%% forced xelatex engine).
\usepackage{amsmath,amssymb,amsthm,mathtools}

\numberwithin{equation}{section}
\newtheorem{thm}{Theorem}[section]
\newtheorem{cor}[thm]{Corollary}
\newtheorem{Lemma}[thm]{Lemma}
\newtheorem{prop}[thm]{Proposition}
 { \theoremstyle{definition}
\newtheorem{defn}[thm]{Definition}
\newtheorem{Note}[thm]{Note}
\newtheorem{Example}[thm]{Example}
\newtheorem{rmk}[thm]{Remark} }
\newcommand{\del}{\partial}
\newcommand{\delb}{\bar{\partial}}\newcommand{\dt}{\frac{\partial}{\partial t}}
\newcommand{\brs}[1]{\left| #1 \right|}
\newcommand{\G}{\Gamma}
\newcommand{\gG}{\Gamma}
\renewcommand{\gg}{\gamma}
\newcommand{\eps}{\epsilon}
\newcommand{\gD}{\Delta}
\newcommand{\gd}{\delta}
\newcommand{\gs}{\sigma}
\newcommand{\gS}{\Sigma}
\newcommand{\gu}{\upsilon}
\newcommand{\gU}{\Upsilon}
\newcommand{\gk}{\kappa}
\newcommand{\gl}{\lambda}
\newcommand{\gt}{\theta}
\newcommand{\gw}{\omega}
\newcommand{\ga}{\alpha}
\newcommand{\gb}{\beta}
\newcommand{\gL}{\Lambda}
\newcommand{\N}{\nabla}
\newcommand{\FF}{\mathcal F}
\newcommand{\WW}{\mathcal W}
\newcommand{\DD}{\mathcal D}
\newcommand{\JJ}{\mathcal J}
\renewcommand{\SS}{\mathcal S}
\newcommand{\EE}{\mathcal E}
\newcommand{\LL}{\mathcal L}
\newcommand{\CC}{\mathcal C}
\newcommand{\R}{\mathbb R}
\newcommand{\til}[1]{\widetilde{#1}}
\newcommand{\nm}[2]{\brs{\brs{ #1}}_{#2}}
\newcommand{\gf}{E}
\newcommand{\boxit}[1]{\fbox{\mbox{$#1$}}}
\newcommand{\ohat}[1]{\overset{\circ}{#1}}
\newcommand{\KN}{{\mathbin{\bigcirc\mkern-15mu\wedge}}}
\renewcommand{\bar}[1]{\overline{#1}}
\newcommand{\dtau}{\frac{d}{d \tau}}
\renewcommand{\i}{\sqrt{-1}}
\renewcommand{\part}{\del}
\newcommand{\bpart}{\bar{\del}}
\newcommand{\bbJ}{\mathbb J}
\newcommand{\ba}{\bar{a}}
\newcommand{\bb}{\bar{b}}
\newcommand{\bc}{\bar{c}}
\newcommand{\bd}{\bar{d}}
\newcommand{\be}{\bar{e}}
\newcommand{\bg}{\bar{g}}
\newcommand{\bh}{\bar{h}}
\newcommand{\bi}{\bar{i}}
\newcommand{\bj}{\bar{j}}
\newcommand{\bk}{\bar{k}}
\newcommand{\bl}{\bar{l}}
\newcommand{\bm}{\bar{m}}
\newcommand{\bn}{\bar{n}}
\newcommand{\bp}{\bar{p}}
\newcommand{\bq}{\bar{q}}
\newcommand{\br}{\bar{r}}
\newcommand{\bs}{\bar{s}}
\newcommand{\bu}{\bar{u}}
\newcommand{\bv}{\bar{v}}
\newcommand{\bw}{\bar{w}}
\newcommand{\bx}{\bar{x}}
\newcommand{\by}{\bar{y}}
\newcommand{\bz}{\bar{z}}
\newcommand{\IP}[1]{\left<#1\right>}
\newcommand{\wnorm}[1]{\left|#1\right|^2_{\frac1k}}
\newcommand{\wwnorm}[1]{\left|#1\right|^2_{\frac{k-1}k}}
\newcommand{\dd}[1]{\frac{\del}{\del #1}}
\newcommand{\bga}{\bar{\alpha}}
\newcommand{\bgb}{\bar{\beta}}
\newcommand{\bgg}{\bar{\gamma}}
\newcommand{\bmu}{\bar{\mu}}
\newcommand{\bnu}{\bar{\nu}}
\newcommand{\bgs}{\bar{\sigma}}
\DeclareMathOperator{\Var}{Var}
\DeclareMathOperator{\USC}{USC}
\DeclareMathOperator{\Sym}{Sym}
\DeclareMathOperator{\Rc}{Rc}
\DeclareMathOperator{\Rm}{Rm}
\DeclareMathOperator{\inj}{inj}
\DeclareMathOperator{\tr}{tr}
\DeclareMathOperator{\Ker}{Ker}
\DeclareMathOperator{\Img}{Im}
\DeclareMathOperator{\Id}{Id}
\DeclareMathOperator{\divg}{div}
\DeclareMathOperator{\grad}{grad}
\DeclareMathOperator{\Vol}{Vol}
\DeclareMathOperator{\diam}{diam}
\DeclareMathOperator{\im}{Im}
\DeclareMathOperator{\Area}{Area}
\DeclareMathOperator{\Ca}{Ca}
%\DeclareMathOperator*{\argmax}{argmax}
%\DeclareMathOperator{\supp}{supp}
\DeclareMathOperator{\Length}{Length}
\DeclareMathOperator{\I}{I}
\DeclareMathOperator{\End}{End}
\DeclareMathOperator{\rank}{Rank}
\DeclareMathOperator{\Ann}{Ann}
\DeclareMathOperator{\Symp}{Symp}
\DeclareMathOperator{\Ham}{Ham}
\DeclareMathOperator{\Diff}{Diff}
\DeclareMathOperator{\Vect}{Vect}
\usepackage{cleveref}
\crefname{originalequation}{original equation}{original equations}
\Crefname{originalequation}{Original equation}{Original equations}
\makeatletter
\newlabel{eq: heatmeasure}{{3.2}{source}{}{sourcepointer-1}{}}
\newlabel{eq: heatmeasure@cref}{{[originalequation][0][]3.2}{[1][0][]source}{}{}{}}
\makeatother
\crefname{originallemma}{original lemma}{original lemmas}
\Crefname{originallemma}{Original lemma}{Original lemmas}
\makeatletter
\newlabel{l:rhoBHE}{{3.5}{source}{}{sourcepointer-2}{}}
\newlabel{l:rhoBHE@cref}{{[originallemma][0][]3.5}{[1][0][]source}{}{}{}}
\makeatother
\crefname{originalproposition}{original proposition}{original propositions}
\Crefname{originalproposition}{Original proposition}{Original propositions}
\makeatletter
\newlabel{prop: energy}{{3.4}{source}{}{sourcepointer-3}{}}
\newlabel{prop: energy@cref}{{[originalproposition][0][]3.4}{[1][0][]source}{}{}{}}
\makeatother
\crefname{originalreference}{original reference}{original references}
\Crefname{originalreference}{Original reference}{Original references}
\makeatletter
\newlabel{thm: existence}{{3.1}{source}{}{sourcepointer-4}{}}
\newlabel{thm: existence@cref}{{[originalreference][0][]3.1}{[1][0][]source}{}{}{}}
\makeatother
\title{Wasserstein distance monotonicity for backward heat flows under generalized Ricci flow}
\author{Kopfer, Eva, Streets, Jeffrey}
\date{doi:10.3842/SIGMA.2024.003; CC BY 4.0}
\begin{document}
\maketitle
\setcounter{section}{3}
\setcounter{equation}{0}
\setcounter{thm}{5}
\noindent\textit{Editorial note.} All delivered mathematical text is the author's own, copied from the pinned journal source file sources/sigma-2024-003-optimal-transport-and-generalized-ricci-flow/pinned-original-source.tex (Optimal Transport and Generalized Ricci Flow); no statement or proof step was rewritten, reordered or removed.\par\medskip
\noindent\textit{Editorial note.} The counted proof is the author's complete argument: the epsilon-near-geodesic characterization of the Wasserstein distance (thm: existence) supplies the initial curve, and the middle bound chains the energy monotonicity of Proposition prop: energy (built card L60069) with the density-form translation of Lemma l:rhoBHE; no proof step is omitted.\par\medskip
\noindent\textit{Editorial note.} The label c:WDM names the counted statement, and the proof environment immediately following it in the source is delivered verbatim as its proof.\par\medskip
\noindent\textit{Editorial note.} The counted statement is carried by the author's own cor environment in the article source.\par\medskip
\noindent\textit{Editorial note.} Journal-unit assembly: card-local preamble (author theorem environments and macros verbatim), article class, no bibliography in the counted span (the span carries no \textbackslash{}cite commands); references leaving the span are declared as navigation stubs carrying the article's printed numbers.\par\medskip
\noindent\textit{Editorial note.} SPAN OVERRIDE: statement=[399, 400] proof=[401, 413] used instead of the recomputed pairing; reason: census pairing defect: the census statement span [399,413] swallows the proof environment that sits inside the author cor environment (begin\{proof\} at 401, end\{proof\} at 412, end\{cor\} at 413), so the census reports pairing none with a foreign next\_proof [443,475]; spans re-verified line by line in the pinned source; the counted proof span carries the closing end\{cor\} line of the embedding environment so the union interval closes (gate closure rule).\par\medskip
\noindent\textit{Editorial scope.} the counted statement and its own proof environment only\par\medskip
\noindent\textit{Allowed standard prerequisites (editorial).}\begin{itemize}
\item The backward heat equation in measure form (eq: heatmeasure). A smooth one-parameter family of probability measures (mu\_t) under generalized Ricci flow satisfies the backward heat flow when d\_t mu = -Delta mu.. Cited at statement line 399: the two solutions mu\textasciicircum{}1\_t, mu\textasciicircum{}2\_t of the backward heat equation are taken in this measure form; declared as a navigation stub carrying the article's printed number.
\item The epsilon-near-geodesic characterization of the Wasserstein distance (thm: existence). For any two probability measures and any eps > 0 there exists a curve mu from mu\textasciicircum{}1 to mu\textasciicircum{}2 whose Lagrangian energy satisfies E(mu) <= 1/2 W\_t(mu\textasciicircum{}1, mu\textasciicircum{}2)\textasciicircum{}2 + eps.. Invoked at proof line 402 to choose the epsilon-near-optimal curve connecting mu\textasciicircum{}1\_\{t0\} and mu\textasciicircum{}2\_\{t0\}; declared as a navigation stub carrying the article's printed number.
\item The backward-heat-equation lemma in density form (l:rhoBHE, printed 3.5). A family mu\_t = rho\_t e\textasciicircum{}\{-f\_t\} dV\_t satisfies the backward heat flow in measure form if and only if the density solves d\_t rho = -Delta\_f rho + R\textasciicircum{}\{H,f\} rho.. Cited at proof line 409 (together with Proposition prop: energy) when the two-parameter density families are used to bound the time-t Wasserstein distance; declared as a navigation stub.
\item Energy monotonicity along the backward heat flow (prop: energy). Along a backward heat flow the Lagrangian energy is nonincreasing in time: E(mu\_t) <= E(mu\_\{t0\}) for t <= t0.. Invoked at proof line 409 to bound E(mu\_t) by E(mu\_\{t0\}) in the middle of the final display; declared as a navigation stub.
\end{itemize}\medskip
\ifdefined\hypertarget\hypertarget{sourcepointer-1}{}\fi
\noindent\textit{Source reference.} Original equation 3.2: sources/sigma-2024-003-optimal-transport-and-generalized-ricci-flow/pinned-original-source.tex, source lines 390--390. Navigation pointer invoked at line(s) 399: the counted text cites that numbered equation of the same article, which is not retained in this card, so the reference is declared as a pointer carrying the number the article itself prints for it (the section and the place of the equation inside that section, derived by the census v3.1 rule for the section-scoped equation counter); the cited equation is an allowed input and none of it is reproduced inside the span.\par\medskip
\ifdefined\hypertarget\hypertarget{sourcepointer-2}{}\fi
\noindent\textit{Source reference.} Original lemma 3.5: sources/sigma-2024-003-optimal-transport-and-generalized-ricci-flow/pinned-original-source.tex, source lines 389--389. Navigation pointer invoked at line(s) 407: the counted statement is phrased in notation from that statement of the same article, which is not retained in this card, so the reference is declared as a pointer carrying the number the article itself prints for it (the section and the place of the statement inside that section, derived by the census v3.1 rule for the shared theorem counter); the cited statement is an allowed input and none of it is reproduced inside the span.\par\medskip
\ifdefined\hypertarget\hypertarget{sourcepointer-3}{}\fi
\noindent\textit{Source reference.} Original proposition 3.4: sources/sigma-2024-003-optimal-transport-and-generalized-ricci-flow/pinned-original-source.tex, source lines 324--324. Navigation pointer invoked at line(s) 407: the counted statement is phrased in notation from that statement of the same article, which is not retained in this card, so the reference is declared as a pointer carrying the number the article itself prints for it (the section and the place of the statement inside that section, derived by the census v3.1 rule for the shared theorem counter); the cited statement is an allowed input and none of it is reproduced inside the span.\par\medskip
\ifdefined\hypertarget\hypertarget{sourcepointer-4}{}\fi
\noindent\textit{Source reference.} Original reference 3.1: sources/sigma-2024-003-optimal-transport-and-generalized-ricci-flow/pinned-original-source.tex, source lines 286--286. Navigation pointer invoked at line(s) 402: the counted statement is phrased in notation from that statement of the same article, which is not retained in this card, so the reference is declared as a pointer carrying the number the article itself prints for it (the section and the place of the statement inside that section, derived by the census v3.1 rule for the shared theorem counter); the cited statement is an allowed input and none of it is reproduced inside the span.\par\medskip
\begin{cor} \label{c:WDM} Let $\big(\mu^1_t\big)$, $\big(\mu^2_t\big)$ be two solutions of the backward heat equation \eqref{eq: heatmeasure}
in $P^\infty(M)$. Then $W_{t}\big(\mu^1_t,\mu^2_t\big)$ is nondecreasing in $t$.

\begin{proof}
 Fix $t_0$. For each $\varepsilon>0$, we may choose according to Theorem \ref{thm: existence} a curve $\mu\colon [0,1]\to P^\infty(M)$ with $\mu(0)=\mu^1_{t_0}$ and $\mu(1)=\mu^2_{t_0}$ satisfying
 \begin{align*}
 E(\mu)\leq \frac12 W_{t_0}\big(\mu^1_{t_0},\mu^2_{t_0}\big)^2+\varepsilon,
 \end{align*}
 where $E(\mu)$ is the Lagrangian of the curve $\mu$ at time $t_0$.
 Let $t\leq t_0$ and let $\mu_t(s)$ be the backward heat flow with $\mu_{t_0}(s)=\mu(s)$. Observe that this implicitly defines two-parameter families $(\rho(s,t), \phi(s,t))$ as described above. Then we know by Proposition \ref{prop: energy} and Lemma \ref{l:rhoBHE} that
 \begin{align*}
 \frac12 W_{t}\big(\mu^1_t,\mu^2_t\big)^2\leq E(\mu_t)\leq E(\mu_{t_0})\leq \frac12 W_{t_0}\big(\mu^1_{t_0},\mu^2_{t_0}\big)^2+\varepsilon.
 \end{align*}
 As $\varepsilon > 0$ is arbitrary, the result follows.
\end{proof}
\end{cor}

\end{document}
